\section{平滑 Scaling Laws 与不平滑任务能力}

前一章确定了证据边界，本章建立最基本的对照。语言模型的 held-out loss 往往随计算、数据和参数呈平滑 power law；但用户真正关心的任务指标可能在很长区间里维持 chance level，然后才快速上升。两条曲线可以同时成立，因为它们测量的对象、尺度和阈值化方式不同。

\subsection{先看论文地图，再看单条曲线}

课件先展示论文首页与任务清单，是为了提醒读者：``emergence'' 不是从一个漂亮案例概括出来的，而是跨多个模型家族、few-shot tasks 和 prompting techniques 汇总得到的现象目录。目录本身不能证明统一机制，却能提出可复用的比较问题。

\lecturefigure{slide-03.jpg}{Emergent Abilities of Large Language Models：论文与能力目录}{官方课件第 3 页；Wei et al., TMLR 2022}

读这张图时先区分“catalog”与“theory”。论文整理了 arithmetic、symbolic manipulation、language understanding 等任务的阈值现象，也整理了某些只在较大模型上有效的 prompting 技术；它并未给出一个能从参数、数据组成直接预测全部阈值的封闭理论。

\subsection{Scaling law：可预测的是平均 loss}

Scaling laws 的典型结论是：当其他条件受控时，test loss 会随模型参数 $N$、训练 token 数 $D$ 或训练计算 $C$ 近似按幂律下降。课堂用 Kaplan et al. 的三组图建立“平滑且可外推”的基线，再把涌现定义为小模型曲线无法直接揭示的任务级变化。这里先理解平均预测误差如何变化，才能避免把后续某个 benchmark 的 abrupt-looking score 误读为与 loss scaling 相矛盾。

\lecturefigure{slide-04.jpg}{模型、数据和计算增长带来的可预测 loss 改进}{官方课件第 4 页；Kaplan et al., 2020}

\begin{knowledgebox}{读图：三张 scaling 曲线在说什么}
横轴分别代表 compute、dataset size 与 parameter count，纵轴是 held-out language-model loss。首先比较 log--log 坐标上的整体直线趋势，而不是某个单点；其次注意这些是分布平均的 next-token prediction 误差，不是某项离散任务的准确率；最后，低 loss 说明模型对数据分布的预测更好，却不直接告诉我们它是否跨过某个任务的 random baseline。
\end{knowledgebox}

训练计算常用一个粗略关系表达：
\[
C \propto N D,
\]
其中 $C$ 是预训练 floating-point operations（FLOPs），$N$ 是参与前向和反向计算的参数规模，$D$ 是训练 token 数。比例常数取决于 Transformer 结构、序列长度、是否重复训练数据以及实现细节，因此这不是精确账本，只是解释为什么固定 $D$ 的模型家族中，参数量与训练 FLOPs 可能给出相似横轴。

\teachervoice{课堂一开始因屏幕共享问题重复了三次同一结论：loss 的改善是可预测的，而本讲所说的 emergence 恰恰是只看小模型时难以预测的任务变化。这段小插曲反而把两类证据的差别说得非常清楚。}

\subsection{从物理直觉到可测定义}

``More is Different'' 的科学直觉帮助理解为什么量变可能带来质变，但语言模型研究仍需要可复验的 operational definition。否则，任何令人惊讶的输出都可以被事后称为涌现，概念就失去比较价值。

\lecturefigure{slide-05.jpg}{科学中的 emergence：定量变化引出定性变化}{官方课件第 5 页；Anderson 1972 与 Jacob Steinhardt 2022}

课堂用两个直觉例子说明“规模改变可达状态”：少量 uranium 与足够密集的 uranium 处在不同反应区域；小分子与 DNA 在信息编码能力上也不是简单线性差别。类比只负责建立直觉，不能替代语言模型中的控制实验，因为参数、数据与算法不像物理质量那样是单一标量。

\lecturefigure{slide-06.jpg}{大语言模型涌现能力的操作性定义与三种 scale 轴}{官方课件第 6 页；Wei et al., 2022}

\begin{importantbox}{操作性定义}
若某项 ability 在较小模型上未被测得、在较大模型上被测得，则称其为 emergent ability。对 few-shot classification，课堂进一步把“未被测得”具体化为接近 random accuracy，把“被测得”具体化为显著超过 random accuracy。这个定义依赖任务、prompt、metric、baseline 与被比较的模型族。
\end{importantbox}

\begin{warningbox}{不要把横轴偷换成“智能”}
Training FLOPs、参数量和训练数据量是三个不同资源轴。即使它们在同一模型家族中高度相关，也不能把其中任意一个直接命名为 intelligence。数据质量、token mixture、architecture、optimizer 与 post-training 都可能改变同一参数量模型的有效能力。
\end{warningbox}

\subsection{本章小结}

平滑 scaling law 描述分布平均的预测误差；emergent-ability plot 描述特定任务在特定 metric 上是否超过 baseline。前者可以连续改进而后者长期不动，因为任务准确率可能要等多个子技能都达到最低可靠度后才显现。

